\subsection{加拉太}
\subsubsection{加拉太範圍-三次宣教都經過加拉太地區}
\begin{table}[H]
\begin{center} 
\begin{tabular}{ p{2.cm}| p{2.5cm}|p{10cm}} 
\toprule 
一次宣教 &13:13-14:25 &建立南加拉太教會（彼西底的安提阿/以哥念/路司得/特庇/别加）\\
二次宣教 &16:1-7&經過南加拉太，堅固諸教會(每西亚/弗吕家)\\
三次宣教 &18:23 &可能經過的是北加拉太(加拉太和弗吕家地方) \\
\bottomrule 
\end{tabular}
\caption{加拉太宣教時間——使徒行傳}
\end{center}
\end{table}

一，相关经文
\begin{itemize} 
\itemsep-.5em 
\item 使徒行传13章-使徒行传18章；
\item 加拉太书；
\item 革勒士往加拉太 (提後4：11）
\end{itemize}

二，保罗第一次宣教旅行在加拉太的事工与经历

\begin{itemize} 
\itemsep-.5em 
\item 保罗在彼西底的安提阿的事工经历\\
保罗和巴拿巴在别加与马可分开之后，就继续上路去了彼西底的安提阿。圣经记载了保罗在那里的服事。
保罗在彼西底的安提阿服事的机会是有管会堂的叫人过去让保罗和巴拿说些劝勉众人的话。保罗才站起来对以色列人开始讲福音。他是从神将以色列出埃及开始讲起，一直讲到主耶稣的死里复活，最后讲述主耶稣就是那位应许的大卫后裔，并鼓励他们相信福音，结果是在那里的人都很喜欢听神的道，许多敬虔的犹太人都跟从了保罗和巴拿巴。
\begin{itemize} 
\itemsep-.5em 
\item 1.	保罗盼望的目的
保罗所盼望的目的是把福音传给住在彼西底的安提阿的人，让他们相信福音（徒13：40），并且劝勉他们恒久的在神的恩中（徒13：43）。他们的目的不只是对犹太人，还有对外邦人传福音（徒13：46）。保罗期望的目的最后应该说达到了，因为在第一个安息日分享之后有很多犹太人和很多敬虔进犹太教的人都跟从了保罗和巴拿巴（徒 13：43）。第二个安息日的时候很多的外邦人也欢喜赞美神的道（徒13：48），主的道也传遍了彼西底的安提阿这一带地方。所以他们满心欢喜，又被圣灵充满（徒13：52）。
\item 2.	保罗所面对的困难
保罗在彼西底的安提阿所面临的困难主要来自某一些犹太人和被犹太人挑唆虔敬尊贵的妇女和城内有名望的人。当人们都喜欢听神的道时，很多人就开始嫉妒他们，驳斥保罗所说的话，并且毁谤（徒13：45），他们不仅自己毁谤还挑唆城里虔敬的妇女和有名望的人逼迫保罗巴拿巴，把他们赶出城。
\begin{itemize} 
\itemsep-.5em 
\item a.	如何胜过
保罗和巴拿巴对他们的回应是放胆的对他们讲要把救恩转向外邦人（13：46），出城的时候对着众人跺脚，表示对他们那些逼迫他们的人很不满，但保罗并没有因为困难而灰心丧气，而是积极的面对困难，并且满心喜乐，被圣灵充满（13：52）。
\item b.	给当地信徒带来的影响
当地的信徒们第一次听到这样的道，对此信息兴趣很大，以至于邀请他们下安息日再讲（13：42）。当地的犹太人和敬虔进犹太教的人很多信从了这道，跟随保罗和巴拿巴。到下个安息日，合城的人几乎都来，这在一个没有现代通讯的时候的影响力是空前的。而且这些人不只是来看热闹，是来听神的道（13：44）。
\end{itemize}
\end{itemize}
\item 保罗在以哥念的事工与经历\\
保罗结束了在彼西底的安提阿事工之后到了坐落于安提阿东南方向的以哥念。他们还是向往常一样进入犹太人会堂讲道。
1.	保罗盼望的目的
保罗在以哥念的目的是在当地的会堂里讲道，使当地的犹太人和希腊人信主（14：1）。他们为了这个目的，在那里住了多日，并且依靠神放胆的讲神国的道，并且行了很多的神迹奇事来帮助他们来达到这个目的。
2.	保罗面临的困难
保罗所面临的困难和在彼西底安提阿的情形相似，有一群不顺从的犹太人迫害保罗，他们耸动外邦人，叫他们心里恼恨弟兄。当跟随犹太人和跟随使徒的人分党之后，外邦人和犹太人的官要凌辱保罗，并且用石头打他们。
3.	如何胜过困难
当保罗和巴拿巴受到不顺从的犹太人攻击的时候他们选择依靠主。在一个很艰难的地方，他们却在那住了多日，依靠主放胆讲道，而且他们行了很多神迹奇事来证明神的道是可信的，从而抵挡攻击。
4.	给当地信徒的结果与影响
保罗和巴拿巴应该是第一个去以哥念传福音的人，许多人第一次听福音时，生命就被改变了，接受了福音，而且这一次还有很多的希腊人接受了福音，这应该是令使徒们很欣慰的事，在当地就撒下了福音的种子。当使徒决定在那里住多日的时候，他们藉着行神迹奇事让人们信心更加坚定，从而使许多人附从他们。福音藉着使徒们的工作在当地产生了不小的影响，以至于惊动了很多的外邦人和犹太人的官长来用石头来攻击使徒保罗，如果影响不大，那些官长肯定不会管使徒们，因为使徒们工作对于他们是一种威胁，所以引来了迫害。可能是跟随使徒的人将消息报告给使徒们，他们才有机会逃出去。这表明当地信仰福音的人和使徒们的关系也有被建立起来，彼此信任。此后以哥念地区的福音工作也是这些附从了使徒的人开展的。
\item 保罗在路司得与特庇及之后返回叙利亚安提阿之前的事工与经历\\
当保罗和巴拿巴逃出以哥念之后就来到路司得。使徒保罗在这里的工作圣灵藉着医治的事件来开展使徒们的事工。
1.	保罗盼望的目的
圣经没有记载这次保罗进会堂讲道，很可能就是在一些人群聚集处讲道，从而让一个生而瘸腿的人听见。通过后面人们的反应，我们可以知道这次听福音的人都是外邦人。使徒保罗这一次的目的是想让外邦人也能听见福音而信主。这次的特殊之处在于保罗专门的是为了给外邦人，让外邦人离弃虚妄，归向创造万有的永生的神。
2.	保罗面临的困难
保罗这次所面临的困难和之前有很多不同。当众人看见保罗所行的神迹时，就把他们当做外邦人的神，宙斯和希耳米。此外还有祭司。牵着牛向使徒献祭。使徒们从没遇到过这种的经历—来传神的福音，却被误认为神。当他们给那些人解释一番之后，有些从安提阿和以哥念来的犹太人出来挑唆众人，众人就用石头打保罗，以至于人们都认为保罗已经死了。
3.	保罗如何胜过
保罗在被称为希米耳之后，他的第一反应是“撕裂”衣服，然后“跳”到众人中间，“喊”着说。这一系列的动词反应出保罗很伤心，并且急迫的想让人们认识这位真神。所以他就把神介绍给路司得人认识。保罗首先声明来传福音的目的。是从神的创造开始讲的，让他们知道神的创造大能，并且将证据显出来。当保罗被石头打后，他战胜的办法就是起来，继续走进城里，准备去下一个地方传福音，并没有被逼迫吓到。
4.	带给当地信徒的结果影响
保罗在路司得首先带来的影响应该是对那个生来瘸腿的人。当瘸腿的人听见福音之后，他的生命完全被福音改变了，因为他对神有了信心。当保罗见他有信心时，随即也医治了他身体上的残疾。圣经对那人的描述是那人就跳起来，而且行走。对当地外邦人的影响是对于神的认识。在听到保罗讲真神之前，他们敬拜假神，给假神祭祀。保罗给他们讲完神之后，他们就不再认为保罗是神了，对神也产生一定的了解。当保罗第二次来探访他们的时候，这里也是有一些弟兄姐妹在。
在特庇的影响圣经讲的篇幅不多，没有讲困难，只讲了保罗去传福音，使好些人做门徒。我们由此可知，很多人接受了福音，并且愿意在真道上成长。最后他们还去了别加讲道。
\end{itemize}

三，保罗第二次宣教旅行在加拉太的事工与经历
保罗第二次宣教去了加拉太，但圣经记载保罗在加拉太地区的事奉没有用很多篇幅。加拉太在当时已经有些教会了（16：2）。
\begin{itemize} 
\itemsep-.5em 
\item 保罗第二次宣教旅行在加拉太地区的目的\\
第二次宣教旅行是保罗提出来的，目的是要去回到从前宣传主道的各城，看望弟兄们境况如何。（15：36）。在41节又说去坚固众教会。加拉太地区是他们第一次宣教之旅的重中之重，所以在第二次中也会是重点。保罗的第二次宣教还有一个重要的任务就是把耶路撒冷使徒和长老所定的条规交给门徒遵守。
\item 保罗在加拉太地区遇到的困难及回应\\
保罗和西拉首先来到特庇，又到了路司得。路司得是从前保罗被石头打的地方，但保罗还是来到这里，而且遇到了提摩太。信主的犹太妇人很可能就是第一次宣教旅行时结下的果子。当保罗想带提摩太去的时候，当地的犹太人一定要给提摩太行割礼，保罗虽然知道这不是必须的，但他并没有与弟兄姐妹发生争执。他将大会的决议给教会宣讲后，教会的信心越发坚固，人数天天增加。接着保罗遇到的第二个困难是圣灵禁止他们在亚细亚讲道，他们选择顺服与否，最后保罗选择顺服圣灵的带领，他们就穿过弗吕家、加拉太到了每西亚边境。当他们想要去庇推尼的时候，圣灵又一次不允许，保罗对这样违背自己意愿的事再次选择顺服圣灵。
\item 第二次宣教旅行在加拉太地区的影响\\
圣经在这些篇幅中记载的不是很多，但我们从作者所写的看出，第二次旅行在加拉太地区是很成功的。因为大家都很乐于接受耶路撒冷大会的决定，并且藉着保罗的劝勉，他们信心越发坚定，人数天天增加。
\end{itemize}

\subsubsection{保罗第三次宣教旅行在加拉太的事工与经历}
当保罗从第二次宣教旅行回来后，第三次宣教旅行又去了加拉太。这次所去的目的和第一次不同，但和第二次有些类似的就是去坚固众门徒。圣经没有记载第三次旅行时在加拉太发生的事。
\subsubsection{加拉太教会属灵情况的进展}
加拉太教会属灵的进展可以说是很大的。因为在第一次宣教旅行时，加拉太人是第一次听福音，虽然当时很多人接受了福音，但对真理仍有许多不明白的，等保罗离开时他们就很快信了别的福音—就如加拉太书里面所记“我稀奇你们这么快离开那藉着基督之恩召你们的，去从别的福音”。他们听从别人，在福音上加上了一些律法，以至于不是福音了（加1：6）。甚至保罗称呼他们为“无知的加拉太人（加3：1）”，可见第一次听福音之后对福音认识是不够的。但当保罗第二次去坚固这些门徒的时候，他们接受了耶路撒冷大会的决定，而且信心开始坚定，人数增加。而且加拉太已经成为一个榜样，他们为圣徒捐钱，被保罗在给哥林多人的信中提到（林前16：1）等到第三次时，加拉太地区应该是非常成熟的地方了。
从加拉太的灵程来看，总体上是呈上升趋势。每一次保罗对他们有教导之后，灵命就有一个稳步的提升，当真理足够让他们自己开始寻求神的时候，他们的灵命就不容易被外界的虚假信息所蒙蔽了双眼。



\subsubsection{加拉太教會的主要問題}
\begin{itemize}
\itemsep-.5em 
\item 這麽快離棄福音，隨從別的福音
\item 靠律法稱義而非信心
\item 守割禮和其他猶太律法的問題
\end{itemize}



\subsubsection{以哥念 Konya （徒13:50-51/14:1-7,19-23，21；提後\textbf{3:1-16}）}
是今日土耳其名叫 Konya 的一個大城，位於安提阿東南約 122 公里，路斯得之北約 30 公里，是當時羅馬帝國呂高尼省的首府，它是在一處地勢略高的平原邊緣，土地肥沃，水源充足，再有帝國的軍事大道通過，故也是一個軍事和商業重鎮，當地的猶太人也很多，該城又以奉拜 Mother Godess 而聞名，其祭司為獨身或係閹人。已有以哥念的錢幣被發現。

\begin{itemize}
\itemsep-.25em 
\item 徒\textbf{13:50-51} 但猶太人挑唆虔敬、尊貴的婦女和城內有名望的人，逼迫保羅、巴拿巴，將他們趕出(彼西底的安提阿)境外。 51 二人對著眾人跺下腳上的塵土，就往\textcolor{red}{以哥念}去了。 52 門徒滿心喜樂，又被聖靈充滿。
\item 徒\textbf{14:1-7} 二人在\textcolor{red}{以哥念}同進猶太人的會堂，在那裡講的叫猶太人和希臘人信的很多。 2 但那不順從的猶太人聳動外邦人，叫他們心裡惱恨弟兄。 3 二人在那裡住了多日，倚靠主放膽講道，主藉他們的手施行神蹟奇事，證明他的恩道。 4 城裡的眾人就分了黨，有附從猶太人的，有附從使徒的。 5 那時，外邦人和猶太人並他們的官長一齊擁上來，要凌辱使徒，用石頭打他們。 6 使徒知道了，就逃往呂高尼的路司得、特庇兩個城和周圍地方去， 7 在那裡傳福音。
\item 徒\textbf{14:19-23} 但有些猶太人從安提阿和\textcolor{red}{以哥念}來，挑唆眾人，就用石頭打保羅，以為他是死了，便拖到城外。 20 門徒正圍著他，他就起來，走進城去。第二天，同巴拿巴往特庇去。 21 對那城裡的人傳了福音，使好些人做門徒，就回路司得、\textcolor{red}{以哥念}、安提阿去， 22 堅固門徒的心，勸他們恆守所信的道，又說：「我們進入神的國，必須經歷許多艱難。」 23 二人在各教會中選立了長老，又禁食禱告，就把他們交託所信的主。 - 14:19-23
\item 徒\textbf{16:1-5} 保羅來到特庇，又到路司得。在那裡有一個門徒名叫提摩太，是信主之猶太婦人的兒子，他父親卻是希臘人。 2 路司得和\textcolor{red}{以哥念}的弟兄都稱讚他。 3 保羅要帶他同去，只因那些地方的猶太人都知道他父親是希臘人，就給他行了割禮。 4 他們經過各城，把耶路撒冷使徒和長老所定的條規交給門徒遵守。 5 於是眾教會信心越發堅固，人數天天加增。
\item 提後\textbf{3:1-16} 你該知道，末世必有危險的日子來到。 2 因為那時人要專顧自己，貪愛錢財，自誇，狂傲，謗讟，違背父母，忘恩負義，心不聖潔， 3 無親情，不解怨，好說讒言，不能自約，性情凶暴，不愛良善， 4 賣主賣友，任意妄為，自高自大，愛宴樂不愛神， 5 有敬虔的外貌卻背了敬虔的實意。這等人你要躲開。 6 那偷進人家，牢籠無知婦女的，正是這等人。這些婦女擔負罪惡，被各樣的私慾引誘， 7 常常學習，終久不能明白真道。 8 從前雅尼和佯庇怎樣敵擋摩西，這等人也怎樣敵擋真道。他們的心地壞了，在真道上是可廢棄的。 9 然而他們不能再這樣敵擋，因為他們的愚昧必在眾人面前顯露出來，像那二人一樣。 10 但你已經服從了我的教訓、品行、志向、信心、寬容、愛心、忍耐， 11 以及我在安提阿、\textcolor{red}{以哥念}、路司得所遭遇的逼迫、苦難。我所忍受是何等的逼迫，但從這一切苦難中，主都把我救出來了。 12 不但如此，凡立志在基督耶穌裡敬虔度日的，也都要受逼迫。 13 只是作惡的和迷惑人的必越久越惡，他欺哄人，也被人欺哄。 14 但你所學習的、所確信的，要存在心裡，因為你知道是跟誰學的， 15 並且知道你是從小明白聖經，這聖經能使你因信基督耶穌有得救的智慧。16 聖經都是神所默示的(凡神所默示的聖經)，於教訓、督責、使人歸正、教導人學義都是有益的， 17 叫屬神的人得以完全，預備行各樣的善事。
\end{itemize}


\subsubsection{路斯得 Lystra (徒14:7-21；16:2-3)}

\begin{itemize}
\itemsep-.25em 
\item 徒\textbf{14:1-23} 二人在\textcolor{red}{路司得}同進猶太人的會堂，在那裡講的叫猶太人和希臘人信的很多。 2 但那不順從的猶太人聳動外邦人，叫他們心裡惱恨弟兄。 3 二人在那裡住了多日，倚靠主放膽講道，主藉他們的手施行神蹟奇事，證明他的恩道。 4 城裡的眾人就分了黨，有附從猶太人的，有附從使徒的。 5 那時，外邦人和猶太人並他們的官長一齊擁上來，要凌辱使徒，用石頭打他們。 6 使徒知道了，就逃往呂高尼的\textcolor{red}{路司得}、特庇兩個城和周圍地方去， 7 在那裡傳福音。8 路司得城裡坐著一個兩腳無力的人，生來是瘸腿的，從來沒有走過。 9 他聽保羅講道，保羅定睛看他，見他有信心，可得痊癒， 10 就大聲說：「你起來，兩腳站直！」那人就跳起來，而且行走。眾人看見保羅所做的事，就用呂高尼的話大聲說：「有神藉著人形降臨在我們中間了！」 12 於是稱巴拿巴為宙斯，稱保羅為希耳米，因為他說話領首。 13 有城外宙斯廟的祭司牽著牛、拿著花圈，來到門前，要同眾人向使徒獻祭。 14 巴拿巴、保羅二使徒聽見，就撕開衣裳，跳進眾人中間，喊著說： 15 「諸君，為什麼做這事呢？我們也是人，性情和你們一樣。我們傳福音給你們，是叫你們離棄這些虛妄，歸向那創造天、地、海和其中萬物的永生神。 16 他在從前的世代，任憑萬國各行其道， 17 然而為自己未嘗不顯出證據來，就如常施恩惠，從天降雨賞賜豐年，叫你們飲食飽足、滿心喜樂。」 18 二人說了這些話，僅僅地攔住眾人不獻祭於他們。但有些猶太人從安提阿和以哥念來，挑唆眾人，就用石頭打保羅，以為他是死了，便拖到城外。 20 門徒正圍著他，他就起來，走進城去。第二天，同巴拿巴往特庇去。 21 對那城裡的人傳了福音，使好些人做門徒，就回路司得、\textcolor{red}{路司得}、安提阿去， 22 堅固門徒的心，勸他們恆守所信的道，又說：「我們進入神的國，必須經歷許多艱難。」 23 二人在各教會中選立了長老，又禁食禱告，就把他們交託所信的主。

\item 徒\textbf{16:1-5} 保羅來到特庇，又到\textcolor{red}{路司得}。在那裡有一個門徒名叫提摩太，是信主之猶太婦人的兒子，他父親卻是希臘人。 2 \textcolor{red}{路司得}和以哥念的弟兄都稱讚他。 3 保羅要帶他同去，只因那些地方的猶太人都知道他父親是希臘人，就給他行了割禮。 4 他們經過各城，把耶路撒冷使徒和長老所定的條規交給門徒遵守。 5 於是眾教會信心越發堅固，人數天天加增。16:2-3
\end{itemize}


\subsubsection{特庇 Derbe  (徒14:21-26；16:1)}

刺 Sting，硝皮匠 Tanner，杜松 Delbeia
●Kerti Huyuk，位於大數以西約 130 公里，路斯得以東約 92 公里，加拉太省的邊境。在當時可能是個相當重要的邊城和軍事據點。



\begin{itemize}
\itemsep-.25em 
\item 對特庇城裡的人傳了福音，使好些人做門徒，就回路司得、以哥念、安提阿去 - 14:20-21
\item 保羅西拉二次宣教第一站就又來到特庇 - 16:1
\item 徒20:4 同他到亞西亞去的、有庇哩亞人畢羅斯的兒子所巴特、帖撒羅尼迦人亞里達古、和西公都、還有特庇人該猶、並提摩太、又有亞西亞人推基古、和特羅非摩。
\end{itemize}


\subsection{旁非利亞/彼西底}

\subsubsection{旁非利亞 - 亞大利 (Antalya)}
亞大魯的城 City of Attalus，和靄的父親 Gentle father
●是今日土耳其南方的一個主要大港埠，大約是在 160 BC 所建，現名 Antalya，位於 Cataractes 河口，距別加約 17 公里，原是別加城的港口，但因日漸繁榮，現已超過並取代了它的地位。

徒14:25 保羅和巴拿巴在別加講道後，就下到亞大利，從那裡坐船往安提阿去。


\subsubsection{旁非利亞 - 別加(Perga)}

要塞 Citadel, 土質的，世俗的 Earthy
●Murtana，位於今日土耳其的南部，Cestrus 河西岸，距海約 13 公里，距他的港埠亞大利(今日的 Antalya 大城)約 17 公里，距安提阿約 155 公里，因有來自彼西底高原的水流灌溉，故土地肥沃，民豐物富，在新約時代，它曾是羅馬帝國旁非利亞省的省城，在現有的遺跡中，尚有一座羅馬式的戲院，估計可容一萬八千餘人，另有華麗的浴池，有頂棚的長街、田徑場等等，都足以表明它曾是一個人口稠密的城市，該城又以奉拜女神 Artemis (Diana)而聞名。

\textbf{徒13:13-14} 保羅和他的同人從帕弗開船，來到旁非利亞的別加，約翰就離開他們，回耶路撒冷去。 14 他們離了別加往前行，來到彼西底的安提阿\\

\textbf{14:24} 回程時再到別加講道後去亞大利。

\begin{itemize}
\itemsep-.25em 
\item 旁非利亞的別加 - 使徒行傳 13:13,14:25，16:4-5
\begin{itemize}
\itemsep-.25em 
\item 約翰就離開保羅和他的同人，回耶路撒冷去 - 13:13
\item 保羅和巴拿巴從別加講道後下到亞大利囘耶路撒冷 - 14:25
\item 保羅西拉二次宣教經過各城，把耶路撒冷使徒和長老定的條規交給門徒遵守。 眾教會信心更堅固，人數天天加增 - 16:4-5
\end{itemize}
\end{itemize}


\subsubsection{彼西底的安提阿 (Yalvach)——徒 13:14-52、14:19-23，21；16:4-5、提後3:1-15}
Yalvach，今日土耳其中部的大城，四面環山，位於 Anthios 河之西岸，以弗所在其西約 350 公里，以哥念在其東南約 120 公里。在弗呂家統治的時代中，就已有了四通八達的公路網，所以是軍事、政治和商業的中心，在 300 BC 前後再經西流基王大肆增建，並改名為安提阿，待軍事大道經過此城之後，就更增加了他的重要性，295 BC 設立彼西底為省，以此城為省會，故又稱為彼西底的安提阿。188 BC 為羅馬帝國所佔，改屬加拉太省，仍以其為省會，直到保羅的時代。

%彼西底的安提阿 - 使徒行傳 13:14-52，21；16:4-5
\begin{itemize}
\itemsep-.25em 
\item 保羅在安息日講道 - 13:14-42
\item 他們出會堂的時候，眾人請他們到下安息日再講這話給他們聽.二人對他們講道，勸他們務要恆久在神的恩中 - 13:43-45
\item 到下安息日，合城的人幾乎都來聚集，要聽神的道。犹太人看见这麽多人，就满心嫉妒，反驳保罗所讲的，并且毁谤他们。挑唆衆人逼迫保羅、巴拿巴，將他們趕出境外 - 13:46-52
\item 對特庇城裡的人傳了福音，使好些人做門徒，就回路司得、以哥念、安提阿去 - 14:21
\item 保羅西拉二次宣教經過各城，把耶路撒冷使徒和長老定的條規交給門徒遵守。 眾教會信心更堅固，人數天天加增 - 16:4-5
\end{itemize}

\textbf{徒13:14-52} 保羅和他的同人從帕弗開船，來到旁非利亞的別加，約翰就離開他們，回耶路撒冷去。 14 他們離了別加往前行，來到彼西底的安提阿，在安息日進會堂坐下。 15 讀完了律法和先知的書，管會堂的叫人過去，對他們說：「二位兄台，若有什麼勸勉眾人的話，請說。」
保羅就站起來，舉手說：「以色列人和一切敬畏神的人，請聽！ 17 這以色列民的神揀選了我們的祖宗，當民寄居埃及的時候抬舉他們，用大能的手領他們出來。 18 又在曠野容忍 a他們約有四十年。 19 既滅了迦南地七族的人，就把那地分給他們為業。 20 此後給他們設立士師，約有四百五十年，直到先知撒母耳的時候。 21 後來他們求一個王，神就將便雅憫支派中基士的兒子掃羅給他們做王四十年。 22 既廢了掃羅，就選立大衛做他們的王，又為他作見證說：『我尋得耶西的兒子大衛，他是合我心意的人，凡事要遵行我的旨意。』 23 從這人的後裔中，神已經照著所應許的為以色列人立了一位救主，就是耶穌。
在他沒有出來以先，約翰向以色列眾民宣講悔改的洗禮。 25 約翰將行盡他的程途，說：『你們以為我是誰？我不是基督。只是有一位在我以後來的，我解他腳上的鞋帶也是不配的。』 26 弟兄們，亞伯拉罕的子孫和你們中間敬畏神的人哪，這救世的道是傳給我們的。 27 耶路撒冷居住的人和他們的官長，因為不認識基督，也不明白每安息日所讀眾先知的書，就把基督定了死罪，正應了先知的預言。 28 雖然查不出他有當死的罪來，還是求彼拉多殺他。 29 既成就了經上指著他所記的一切話，就把他從木頭上取下來，放在墳墓裡。 30 神卻叫他從死裡復活。 31 那從加利利同他上耶路撒冷的人多日看見他，這些人如今在民間是他的見證。 32 我們也報好信息給你們，就是：那應許祖宗的話， 33 神已經向我們這做兒女的應驗，叫耶穌復活了。正如《詩篇》第二篇上記著說：『你是我的兒子，我今日生你。』 34 論到神叫他從死裡復活，不再歸於朽壞，就這樣說：『我必將所應許大衛那聖潔、可靠的恩典賜給你們。』 35 又有一篇上說：『你必不叫你的聖者見朽壞。』 36 大衛在世的時候遵行了神的旨意，就睡了 b，歸到他祖宗那裡，已見朽壞。 37 唯獨神所復活的，他並未見朽壞。
所以弟兄們，你們當曉得：赦罪的道是由這人傳給你們的！ 39 你們靠摩西的律法，在一切不得稱義的事上信靠這人，就都得稱義了。 40 所以你們務要小心，免得先知書上所說的臨到你們。 41 主說：『你們這輕慢的人，要觀看，要驚奇，要滅亡！因為在你們的時候，我行一件事，雖有人告訴你們，你們總是不信。』」
42 他們出會堂的時候，眾人請他們到下安息日再講這話給他們聽。 43 散會以後，猶太人和敬虔進猶太教的人多有跟從保羅、巴拿巴的。二人對他們講道，勸他們務要恆久在神的恩中。
44 到下安息日，合城的人幾乎都來聚集，要聽神的道。 45 但猶太人看見人這樣多，就滿心嫉妒，硬駁保羅所說的話，並且毀謗。 46 保羅和巴拿巴放膽說：「神的道先講給你們原是應當的，只因你們棄絕這道，斷定自己不配得永生，我們就轉向外邦人去。 47 因為主曾這樣吩咐我們說：『我已經立你做外邦人的光，叫你施行救恩，直到地極。』」 48 外邦人聽見這話，就歡喜了，讚美神的道，凡預定得永生的人都信了。 49 於是主的道傳遍了那一帶地方。但猶太人挑唆虔敬、尊貴的婦女和城內有名望的人，逼迫保羅、巴拿巴，將他們趕出境外。 51 二人對著眾人跺下腳上的塵土，就往以哥念去了。 52 門徒滿心喜樂，又被聖靈充滿。\\


\textbf{徒14:19-23} 但有些猶太人從安提阿和以哥念來，挑唆眾人，就用石頭打保羅，以為他是死了，便拖到城外。 20 門徒正圍著他，他就起來，走進城去。第二天，同巴拿巴往特庇去。 21 對那城裡的人傳了福音，使好些人做門徒，就回路司得、以哥念、安提阿去， 22 堅固門徒的心，勸他們恆守所信的道，又說：「我們進入神的國，必須經歷許多艱難。」 23 二人在各教會中選立了長老，又禁食禱告，就把他們交託所信的主。\\

\textbf{提後3:1-15} 你該知道，末世必有危險的日子來到。 2 因為那時人要專顧自己，貪愛錢財，自誇，狂傲，謗讟，違背父母，忘恩負義，心不聖潔， 3 無親情，不解怨，好說讒言，不能自約，性情凶暴，不愛良善， 4 賣主賣友，任意妄為，自高自大，愛宴樂不愛神， 5 有敬虔的外貌卻背了敬虔的實意。這等人你要躲開。 6 那偷進人家，牢籠無知婦女的，正是這等人。這些婦女擔負罪惡，被各樣的私慾引誘， 7 常常學習，終久不能明白真道。 8 從前雅尼和佯庇怎樣敵擋摩西，這等人也怎樣敵擋真道。他們的心地壞了，在真道上是可廢棄的。 9 然而他們不能再這樣敵擋，因為他們的愚昧必在眾人面前顯露出來，像那二人一樣。 10 但你已經服從了我的教訓、品行、志向、信心、寬容、愛心、忍耐， 11 以及我在安提阿、以哥念、路司得所遭遇的逼迫、苦難。我所忍受是何等的逼迫，但從這一切苦難中，主都把我救出來了。 12 不但如此，凡立志在基督耶穌裡敬虔度日的，也都要受逼迫。 13 只是作惡的和迷惑人的必越久越惡，他欺哄人，也被人欺哄。 14 但你所學習的、所確信的，要存在心裡，因為你知道是跟誰學的， 15 並且知道你是從小明白聖經，這聖經能使你因信基督耶穌有得救的智慧。

